\chapter{Additional Results}
\label{ch:appendix}

\section{Highest-uncertainty cities}

\Cref{tab:appendix-cities} lists the twenty test-region cities with the highest estimated difficulty $\hat\sigma(x)$, with the local observation-density diagnostics mentioned in \cref{sec:geography}, present the number of other cities within $250\,\mathrm{km}$, $500\,\mathrm{km}$ and $1{,}000\,\mathrm{km}$, and the great-circle distance required to reach the $20$th-nearest city.
\begin{table}[H]
\centering
\caption{Twenty highest-uncertainty cities, with local observation-density diagnostics.}
\label{tab:appendix-cities}
\scriptsize
\begin{tabular}{lccccc}
\toprule
City (Country) & $\hat\sigma$ & Within 250 km & Within 500 km & Within 1000 km & Dist.\ to 20th NN (km) \\
\midrule
La Paz (Bolivia) & 5.30 & 2 & 11 & 21 & 948.7 \\
Tacna (Peru) & 5.27 & 2 & 7 & 22 & 952.4 \\
Arica (Chile) & 5.27 & 2 & 7 & 22 & 952.4 \\
Puno (Peru) & 5.25 & 2 & 9 & 22 & 894.0 \\
Juliaca (Peru) & 5.25 & 2 & 9 & 22 & 894.0 \\
Arequipa (Peru) & 5.25 & 4 & 10 & 23 & 839.4 \\
Cochabamba (Bolivia) & 5.25 & 3 & 9 & 21 & 945.1 \\
Cusco (Peru) & 4.95 & 4 & 8 & 23 & 828.1 \\
P\^orto Velho (Brazil) & 4.90 & 0 & 2 & 4 & 1307.6 \\
Rio Branco (Brazil) & 4.85 & 1 & 2 & 19 & 1029.3 \\
Ayacucho (Peru) & 4.66 & 3 & 11 & 23 & 808.0 \\
Ica (Peru) & 4.52 & 4 & 11 & 23 & 903.2 \\
Chincha Alta (Peru) & 4.40 & 5 & 7 & 21 & 991.2 \\
Antofagasta (Chile) & 4.35 & 1 & 2 & 23 & 834.0 \\
Ji Paran\'a (Brazil) & 4.28 & 1 & 2 & 8 & 1366.2 \\
Iquique (Chile) & 4.25 & 2 & 10 & 21 & 907.0 \\
Calama (Chile) & 4.20 & 1 & 6 & 23 & 780.3 \\
Huancayo (Peru) & 4.17 & 5 & 10 & 24 & 907.3 \\
Oruro (Bolivia) & 4.16 & 4 & 11 & 20 & 966.5 \\
Chosica (Peru) & 3.84 & 3 & 7 & 22 & 981.3 \\
\bottomrule
\end{tabular}
\end{table}

\section{Software}

All modelling, calibration and figure generation was implemented in Python by using \texttt{NumPy} and \texttt{SciPy} for the core spherical-harmonic and associated-Legendre computations, whereas \texttt{PyTorch} used for GPU-accelerated design-matrix assembly and linear solves during hyperparameter search and grid reconstruction. In addition, \texttt{pandas} and \texttt{scikit-learn} were used for data handling and cross-validation splitting, while \texttt{Matplotlib}/\texttt{seaborn} was used for visualisation.

All coding used to produce the results and figures in this dissertation was carried out in Python. The full code is available at the following GitHub repository: \url{https://github.com/szeying-code/temperature-scalar-field-reconstruction-with-uncertainty-}. 

\section{Reproducibility details}

For reproducing details, the following settings specify the main fixed parameters used in the pipeline reported in \cref{ch:results}. A random seed of $42$ is used for all data-splitting operations, at the four-way train/validation/calibration/test split of \cref{ch:data}. The five-fold cross-validation folds used both for hyperparameter search and for the difficulty-estimator residual field, with target coverage level of $\alpha=0.1$ throughout, and a difficulty-estimator neighbourhood size of $k_\sigma=20$. A numerical floor of $\epsilon=10^{-6}$ on $\hat\sigma(x)$ applied to avoid division by zero in the normalised nonconformity score. The finite-sample-corrected conformal quantile level $\lceil(1-\alpha)(n_{cal}+1)\rceil/n_{cal}$ with the conservative (\texttt{"higher"}) interpolation rule implemented for empirical quantile estimation.

The calibration threshold obtained with $L=26$ model in \cref{sec:main-results} was $\hat q_\alpha = 2.389$, and the adaptive interval half-width at any query point was then $\hat q_\alpha\hat\sigma(x)$. This threshold of $\hat q_\alpha = 2.389$ is then used with the local difficulty estimate $\hat\sigma(x)$ to determine the interval width. This fixed model prediction intervals can be reconstructed using the calibration threshold $\hat q_\alpha = 2.389$, with the fitted coefficient vector $v$ and the stored training-set out-of-fold residual field used to calculate $\hat\sigma(x)$. Thus, it will be unnecessary to repeat the calibration procedure to make the specific numerical results in \cref{ch:results} reproducible. The accompanying code are able to use the same fixed values to reproduce the reported prediction intervals and results. This prevent changes caused by stochastic variation in cross-validation fold assignment or grid-search order.

